\hspace{3mm}

\centerline{\bf \Huge Домашнее задание.}
\centerline{\bf \Huge Числовые таблицы и процентное содержание.}

\begin{flushright}
    \scriptsize{Если я сдамся из-за того, что испугался,\\
                то не смогу стать сильнее.
                \\Сон Джин-Ву. Поднятие уровня в одиночку}
\end{flushright}


\problem{1!} Литровый напиток Фёдора в два раза слаще полулитрового напитка Баты. Сколько мл сахара в напитке Баты, если в напитке Фёдора концентрация сахара 30\%?

\problem{2!} На шахматной доске произвольного размера в каждой чёрной клетке написали четное число, а в каждой белой $-$ нечётное. При каких размерах доски сумма чисел в каждом столбце будет четной?

\problem{3} В таблице $12\times12$ расставлены различные натуральные числа, сумма которых равна 2S. Известно, что в каждой строке числа возрастают слева направо, а в каждом столбце $-$ снизу вверх. Может ли сумма чисел в центральном квадрате $6\times6$ быть больше S?

\problem{4} Дамир и Вова $-$ отличные высушиватели рыбы (рыба состоит из сухого вещества и воды). За час сушки из одного килограмма рыбы Дамира испарилась половина воды, а за это же время у Вовы процентное содержание сухого вещества выросло в два раза в одном килограмме такой же рыбы. Кто из мальчиков более эффективно сушит рыбу, если минимальное количество воды в рыбе до сушки 1\%, а максимальное $-$ 99\%?